% ====================================================================
% CAND-FO3-1 : 합성 관측 dQ/dV = C_bg + sum_j Q_j xi_eq(1-xi_eq)/w_j  (N9, eq:sum)
%   배치 제안: §10 (sec:sum) 식~(eq:sum) 박스 직후. fig:staging(분리 peak)의 짝 그림.
%   좌표 = 식 그대로의 수치 평가. fig:staging 과 동일한 tab:staging 초기값
%     (U={0.210,0.140,0.120,0.085} V, w={0.020,0.016,0.014,0.012} V, Q={0.10,0.12,0.25,0.50} Q_cell),
%     방전 sigma_d=+1, 배경 C_bg=0.6 Q_cell/V(예시 상수). python 평가·하드코딩.
%   사용 라이브러리: arrows.meta (ch1 preamble 에 로드됨).
% ====================================================================
\begin{figure}[t]
\centering
\begin{tikzpicture}[x=38cm,y=0.36cm,>=Latex]
% ---- axes ----
\draw[-{Latex[length=1.6mm]}] (0.028,0) -- (0.268,0) node[below,font=\scriptsize] {$V_n$ [V]};
\draw[-{Latex[length=1.4mm]}] (0.028,0) -- (0.028,13.4) node[above,font=\scriptsize,align=center] {$\dd Q/\dd V$\\{\tiny[$Q_\cell/$V]}};
\foreach \vv in {0.05,0.10,0.15,0.20,0.25} {\draw (\vv,0.18)--(\vv,-0.18) node[below,font=\scriptsize] {\vv};}
\foreach \yy in {4,8,12} {\draw (0.0268,\yy)--(0.0292,\yy) node[left,font=\scriptsize] {\yy};}
% ---- background pedestal C_bg (dotted horizontal) ----
\draw[densely dotted] (0.030,0.6) -- (0.260,0.6);
\node[font=\tiny,anchor=west] at (0.212,0.42) {$C_\bg$ (예시 상수)};
% ---- four individual transition contributions (thin gray, +C_bg baseline) ----
\draw[black!45,thin] plot[smooth] coordinates {(0.1100,0.6332) (0.1152,0.6429) (0.1203,0.6553) (0.1255,0.6711) (0.1307,0.6913) (0.1359,0.7170) (0.1410,0.7494) (0.1462,0.7900) (0.1514,0.8404) (0.1566,0.9022) (0.1617,0.9769) (0.1669,1.0653) (0.1721,1.1673) (0.1772,1.2813) (0.1824,1.4034) (0.1876,1.5271) (0.1928,1.6438) (0.1979,1.7428) (0.2031,1.8136) (0.2083,1.8477) (0.2134,1.8408) (0.2186,1.7937) (0.2238,1.7124) (0.2290,1.6063) (0.2341,1.4862) (0.2393,1.3621) (0.2445,1.2422) (0.2497,1.1319) (0.2548,1.0343) (0.2600,0.9505)};
\draw[black!45,thin] plot[smooth] coordinates {(0.0600,0.6499) (0.0655,0.6700) (0.0710,0.6981) (0.0766,0.7369) (0.0821,0.7904) (0.0876,0.8631) (0.0931,0.9606) (0.0986,1.0885) (0.1041,1.2515) (0.1097,1.4510) (0.1152,1.6820) (0.1207,1.9293) (0.1262,2.1657) (0.1317,2.3550) (0.1372,2.4611) (0.1428,2.4611) (0.1483,2.3550) (0.1538,2.1657) (0.1593,1.9293) (0.1648,1.6820) (0.1703,1.4510) (0.1759,1.2515) (0.1814,1.0885) (0.1869,0.9606) (0.1924,0.8631) (0.1979,0.7904) (0.2034,0.7369) (0.2090,0.6981) (0.2145,0.6700) (0.2200,0.6499)};
\draw[black!45,thin] plot[smooth] coordinates {(0.0500,0.7187) (0.0548,0.7667) (0.0597,0.8335) (0.0645,0.9261) (0.0693,1.0533) (0.0741,1.2265) (0.0790,1.4585) (0.0838,1.7630) (0.0886,2.1511) (0.0934,2.6262) (0.0983,3.1762) (0.1031,3.7649) (0.1079,4.3277) (0.1128,4.7785) (0.1176,5.0313) (0.1224,5.0313) (0.1272,4.7785) (0.1321,4.3277) (0.1369,3.7649) (0.1417,3.1762) (0.1466,2.6262) (0.1514,2.1511) (0.1562,1.7630) (0.1610,1.4585) (0.1659,1.2265) (0.1707,1.0533) (0.1755,0.9261) (0.1803,0.8335) (0.1852,0.7667) (0.1900,0.7187)};
\draw[black!45,thin] plot[smooth] coordinates {(0.0300,1.0173) (0.0340,1.1761) (0.0379,1.3930) (0.0419,1.6870) (0.0459,2.0814) (0.0498,2.6032) (0.0538,3.2803) (0.0578,4.1359) (0.0617,5.1787) (0.0657,6.3880) (0.0697,7.6977) (0.0736,8.9859) (0.0776,10.0826) (0.0816,10.8046) (0.0855,11.0118) (0.0895,10.6615) (0.0934,9.8255) (0.0974,8.6608) (0.1014,7.3521) (0.1053,6.0592) (0.1093,4.8891) (0.1133,3.8946) (0.1172,3.0872) (0.1212,2.4532) (0.1252,1.9674) (0.1291,1.6016) (0.1331,1.3299) (0.1371,1.1298) (0.1410,0.9835) (0.1450,0.8770)};
% ---- composite total (thick) ----
\draw[very thick] plot[smooth] coordinates {(0.0300,1.0544) (0.0340,1.2284) (0.0381,1.4666) (0.0421,1.7905) (0.0461,2.2264) (0.0502,2.8050) (0.0542,3.5582) (0.0582,4.5129) (0.0623,5.6799) (0.0663,7.0376) (0.0704,8.5146) (0.0744,9.9809) (0.0784,11.2606) (0.0825,12.1745) (0.0865,12.6029) (0.0905,12.5361) (0.0946,12.0777) (0.0986,11.3989) (0.1026,10.6696) (0.1067,10.0060) (0.1107,9.4520) (0.1147,8.9908) (0.1188,8.5730) (0.1228,8.1460) (0.1268,7.6747) (0.1309,7.1479) (0.1349,6.5750) (0.1389,5.9767) (0.1430,5.3765) (0.1470,4.7962) (0.1511,4.2535) (0.1551,3.7621) (0.1591,3.3312) (0.1632,2.9661) (0.1672,2.6682) (0.1712,2.4352) (0.1753,2.2622) (0.1793,2.1420) (0.1833,2.0659) (0.1874,2.0242) (0.1914,2.0068) (0.1954,2.0036) (0.1995,2.0048) (0.2035,2.0018) (0.2075,1.9880) (0.2116,1.9585) (0.2156,1.9111) (0.2196,1.8460) (0.2237,1.7654) (0.2277,1.6729) (0.2318,1.5730) (0.2358,1.4700) (0.2398,1.3680) (0.2439,1.2702) (0.2479,1.1789) (0.2519,1.0956) (0.2560,1.0211) (0.2600,0.9555)};
% ---- transition labels (at each individual peak) ----
\node[font=\tiny,black!55,anchor=south east] at (0.0800,11.15) {$2\!\to\!1$};
\node[font=\tiny,black!55,anchor=west] at (0.1230,5.10) {$2\mathrm L\!\to\!2$};
\node[font=\tiny,black!55,anchor=south west] at (0.1400,2.50) {$3\!\to\!2\mathrm L$};
\node[font=\tiny,black!55,anchor=south] at (0.2083,1.90) {$4\!\to\!3$};
% ---- overlap annotation ----
\draw[->,gray] (0.1150,9.0) -- (0.1010,10.3);
\node[font=\tiny,gray,anchor=west,align=left] at (0.1150,8.6) {겹침 shoulder\\($2\!\to\!1{+}2\mathrm L\!\to\!2$)};
\node[font=\scriptsize,anchor=west] at (0.1120,12.20) {합성(관측) $=$ 굵은 실선};
% ---- discharge direction ----
\draw[->,semithick] (0.150,12.6) -- (0.230,12.6);
\node[font=\tiny,anchor=south] at (0.190,12.6) {방전(탈리튬화) $\xi\!\uparrow$};
\end{tikzpicture}
\caption{합산 노드(N9)가 내는 \emph{관측} $\dd Q/\dd V$ 곡선 --- 배경과 네 전이 봉우리의 선형 합(식~\eqref{eq:sum};
좌표는 식 그대로의 수치 평가). 얇은 회색은 전이별 기여 $Q_j\,\xi_\eq(1-\xi_\eq)/w_j$(식~\eqref{eq:eqpeak}),
굵은 실선은 그 합에 배경을 얹은 합성 곡선이며, 이 굵은 실선이 실제 측정되는 한 곡선이다. 입력은
그림~\ref{fig:staging} 과 같은 표~\ref{tab:staging} 초기값($U_j{=}0.210/0.140/0.120/0.085$ V,
$w_j{=}0.020/0.016/0.014/0.012$ V, $Q_j{=}0.10/0.12/0.25/0.50\,Q_\cell$, 방전 $\sigma_d{=}{+}1$)이고,
배경 $C_\bg$ 는 예시 상수 pedestal($0.6\,Q_\cell/$V --- 모델 입력으로 상수 또는 $V$ 함수)이다. 그림~\ref{fig:staging}
이 네 봉우리를 \emph{분리}해 보였다면, 이 그림은 그것들이 배경 위에서 \emph{합쳐진} 모습을 보인다 --- 중심이
가까운 $2\!\to\!1\cdot2\mathrm L\!\to\!2$ 가 한 봉우리로 겹쳐(shoulder) 서론의 ``peak 겹침'' 관측을 그대로
드러내고, 멀리 떨어진 $4\!\to\!3$ 만 별개 봉우리로 남는다. 순높이는 전이마다 $Q_j/(4w_j)$ 로 크게 달라
($2\!\to\!1$ 이 지배) 합성 곡선의 모양을 $2\!\to\!1$ 봉우리가 이끈다.}
\label{fig:cand-fo3-1}
\end{figure}
